\documentclass[11pt,a4paper]{article}
\usepackage[margin=0.85in]{geometry}
\usepackage[T1]{fontenc}
\usepackage{lmodern}
\usepackage{amsmath,amssymb,booktabs,tabularx,enumitem}
\usepackage[hidelinks]{hyperref}
\setlist[itemize]{nosep,leftmargin=*}
\setlength{\parindent}{0pt}
\setlength{\parskip}{5pt}
\title{\vspace{-1.6cm}Midterm Research Report\\[3pt]\large From Cayley Graphs to Face-Glued Three-Dimensional Spaces}
\author{}
\date{October 30, 2026}
\begin{document}
\maketitle
\vspace{-1.2cm}

\section*{1. Learning and Research Development}
The following summarizes my recent reading and discussions, building on earlier project materials.
\begin{itemize}
\item \textbf{Group theory and representation theory:} Constructed Cayley graphs and studied their spectra through irreducible representations.
\item \textbf{Graph embeddings and surfaces:} Studied face systems, rotation systems, Euler characteristic, and genus.
\item \textbf{Geometry and equivalence:} Distinguished abstract embeddings, spatial realizations, map isomorphism, homeomorphism, and homotopy equivalence.
\item \textbf{Algebraic topology:} Introduced boundary matrices, homology, vertex links, and Hodge Laplacians as tools for checking and comparing constructions.
\item \textbf{Three-dimensional constructions:} Began studying tetrahedral face gluings, boundary configurations, and algorithms parameterized by treewidth.
\end{itemize}

\section*{2. Main Insight and Current Question}
A graph does not determine which cycles are filled as faces. Fixing a graph therefore fixes its ordinary graph Laplacian, but need not fix the topology of an associated space.

\textbf{Current question.} For a fixed tetrahedral face-pairing graph $\Gamma$, can different gluing maps produce closed, connected, orientable $3$-manifolds with isomorphic integral homology but different homeomorphism types?

\textbf{Related question.} Which topological information is lost when different partial constructions are represented by the same boundary configuration?

\section*{3. Current Progress}
\begin{itemize}
\item Reviewed the main ideas of Cayley spectra, genus algorithms, spatial graph embeddings, and boundary-state methods.
\item Rechecked an existing catalogue of $25$ surface examples; all passed the current validator, with numerical Hodge-nullity checks.
\item Worked through Example~7 of Burton--Pettersson: two vertex links are disks and the pinched vertex link is an annulus.
\item Identified an apparent inconsistency in its printed gluing formula; the reconstruction follows the accompanying figure and configuration tables.
\item Complete three-dimensional enumeration and a certified same-homology, non-homeomorphic pair remain unfinished.
\end{itemize}

\section*{4. Initial Scope and Working Hypotheses}
Start with one tetrahedron ($36$ labeled gluing assignments), then a fixed two-tetrahedron graph. The primary graph family has doubled edges along a path and one loop at each endpoint; its treewidth is $1$ for at least two vertices.
\begin{itemize}
\item \textbf{H1:} Some fixed graph in this family supports same-homology, non-homeomorphic manifolds.
\item \textbf{H2:} Such a pair may differ in only one face-gluing map.
\end{itemize}
These are working hypotheses, not established results or claims of previously unknown conjectures.

\newpage
\section*{5. Seven-Week Plan}
The schedule assumes seven weeks beginning October~31; dates can be adjusted to the actual submission deadline.

\renewcommand{\arraystretch}{1.35}
\begin{tabularx}{\textwidth}{@{}p{2.8cm}X@{}}
\toprule
\textbf{Week} & \textbf{Main task}\\
\midrule
1: Oct.\ 31--Nov.\ 6 & Finalize the midterm report, fix the research scope, and independently reconstruct Example~7.\\
2: Nov.\ 7--13 & Build and validate the tetrahedral gluing and vertex-link representation.\\
3: Nov.\ 14--20 & Enumerate the $36$ one-tetrahedron assignments and compute integral homology of valid manifolds.\\
4: Nov.\ 21--27 & Extend to a fixed two-tetrahedron graph and record boundary-state histories.\\
5: Nov.\ 28--Dec.\ 4 & Analyze same-homology candidates and distinguish proved conclusions from unresolved pairs.\\
6: Dec.\ 5--11 & Complete the final report draft, figures, and evidence tables.\\
7: Dec.\ 12--18 & Reproduce the core results, revise the report, and prepare the presentation.\\
\bottomrule
\end{tabularx}

\section*{6. Expected Outcome}
A small, reproducible study connecting surface examples, vertex-link validation, integral homology, and boundary-state information. A new classification theorem or a successful counterexample is not required for completion.

Known results, finite computations, numerical checks, and unproved hypotheses will be stated separately. Different triangulation signatures alone will not be treated as proof of non-homeomorphism. Larger searches and spectral clustering experiments are optional extensions.

\begin{thebibliography}{9}
\small
\bibitem{liu}
X. Liu and S. Zhou,
\emph{Eigenvalues of Cayley graphs},
\href{https://arxiv.org/abs/1809.09829}{arXiv:1809.09829}.

\bibitem{brinkmann}
G. Brinkmann,
\emph{A Practical Algorithm for the Computation of the Genus},
\href{https://arxiv.org/abs/2005.08243}{arXiv:2005.08243}.

\bibitem{barthel}
S. Barthel and F. Buccoliero,
\emph{Constructing embedded surfaces for cellular embeddings of leveled spatial graphs},
\href{https://arxiv.org/abs/2406.03800}{arXiv:2406.03800}.

\bibitem{burton}
B. A. Burton and W. Pettersson,
\emph{Fixed Parameter Tractable Algorithms in Combinatorial Topology},
\href{https://arxiv.org/abs/1402.3876v2}{arXiv:1402.3876v2}.
\end{thebibliography}
\end{document}
